% GENERATED FILE. Edit canonical lessons, then run npm run generate:books.
% canonical-source-hash: fnv1a64:8e02117d47f188cd
% canonical-lessons: GU-C229-kadu, GU-C229-angutho, GU-C229-ghuti, GU-C229-edi, GU-C229-pasli

\chapter{From the wrist to the heel}
\label{ch:gu-from-the-wrist-to-the-heel}
\hlchaptermodality{229}

\begin{chapteropening}
\textbf{By the end of this chapter:} \emph{I can say which part of my hand or foot hurts.}

Complete the last lesson of chapter 229: I can say which part of my hand or foot hurts.
\end{chapteropening}

\section[\texorpdfstring{\gu{કાંડું} (kā̃ḍũ)}{કાંડું (kā̃ḍũ)}]{\gu{કાંડું} (kā̃ḍũ) — wrist}

\label{lesson:GU-C229-kadu}



\begin{quote}
Before the new one: say the Gujarati for skin, then the Gujarati for blister.
\end{quote}

\subsection*{You'll want to know: \gu{કાંડું}}
\textbf{\gu{કાંડું}} — \emph{kā̃ḍũ} — \textquotedblleft{}wrist\textquotedblright{}.

\begin{quote}
\textbf{\gu{શું} \gu{તમારું} \gu{કાંડું} \gu{દુખે} \gu{છે}?} — \emph{śũ tamārũ kā̃ḍũ dukhe che?} — \textquotedblleft{}Does your wrist hurt?\textquotedblright{}
\end{quote}

\begin{cousinweb}
\textbf{\gu{કાંડું}} (\emph{kā̃ḍũ}) is a tadbhava. CDIAL lists it, glossed \textquotedblleft{}wrist\textquotedblright{}, under Sanskrit \emph{kāṇḍa}, \textquotedblleft{}single joint of a plant\textquotedblright{}, a word it thinks probably came from Dravidian, so the wrist is a joint (CDIAL 3023).
\end{cousinweb}

\begin{grammarlens}[title={What you've built}]
One more word for everyday life at A2.
\end{grammarlens}

\subsection*{Your turn}
\begin{itemize}
\raggedright
  \item \emph{Say it:} \emph{kā̃ḍũ}
  \item \emph{Say it:} \emph{kā̃ḍũ}, once more
  \item \emph{Say it:} \emph{śũ tamārũ kā̃ḍũ dukhe che?} — Does your wrist hurt?
  \item \emph{Recall:} say the Gujarati for skin, then the Gujarati for blister, then say \emph{kā̃ḍũ} again
\end{itemize}

\begin{tcolorbox}[breakable,colback=teal!4,colframe=teal!35!black,title={Before you move on}]
What does \gu{કાંડું} mean? (\textquotedblleft{}Wrist\textquotedblright{}.) Say it once more.
\end{tcolorbox}
\section[\texorpdfstring{\gu{અંગૂઠો} (angūṭho)}{અંગૂઠો (angūṭho)}]{\gu{અંગૂઠો} (angūṭho) — thumb, big toe}

\label{lesson:GU-C229-angutho}



\begin{quote}
Before the new one: say the Gujarati for blister, then the Gujarati for wrist.
\end{quote}

\subsection*{You'll want to know: \gu{અંગૂઠો}}
\textbf{\gu{અંગૂઠો}} — \emph{angūṭho} — \textquotedblleft{}thumb, big toe\textquotedblright{}.

\begin{quote}
\textbf{\gu{મારો} \gu{અંગૂઠો} \gu{નથી} \gu{દુખતો}.} — \emph{māro angūṭho nathī dukhto} — \textquotedblleft{}My thumb does not hurt.\textquotedblright{}
\end{quote}

\begin{cousinweb}
\textbf{\gu{અંગૂઠો}} (\emph{angūṭho}) is a tadbhava, a word inherited from Sanskrit and reshaped by sound change. It comes from \emph{aṅguṣṭha}, \textquotedblleft{}thumb; big toe\textquotedblright{} (CDIAL 137).
\end{cousinweb}

\begin{grammarlens}[title={What you've built}]
One more word for everyday life at A2.
\end{grammarlens}

\subsection*{Your turn}
\begin{itemize}
\raggedright
  \item \emph{Say it:} \emph{angūṭho}
  \item \emph{Say it:} \emph{angūṭho}, once more
  \item \emph{Say it:} \emph{māro angūṭho nathī dukhto} — My thumb does not hurt.
  \item \emph{Recall:} say the Gujarati for blister, then the Gujarati for wrist, then say \emph{angūṭho} again
\end{itemize}

\begin{tcolorbox}[breakable,colback=teal!4,colframe=teal!35!black,title={Before you move on}]
What does \gu{અંગૂઠો} mean? (\textquotedblleft{}Thumb, big toe\textquotedblright{}.) Say it once more.
\end{tcolorbox}
\section[\texorpdfstring{\gu{ઘૂંટી} (ghū̃ṭī)}{ઘૂંટી (ghū̃ṭī)}]{\gu{ઘૂંટી} (ghū̃ṭī) — ankle}

\label{lesson:GU-C229-ghuti}



\begin{quote}
Before the new one: say the Gujarati for wrist, then the Gujarati for thumb.
\end{quote}

\subsection*{You'll want to know: \gu{ઘૂંટી}}
\textbf{\gu{ઘૂંટી}} — \emph{ghū̃ṭī} — \textquotedblleft{}ankle\textquotedblright{}.

\begin{quote}
\textbf{\gu{શું} \gu{તમારી} \gu{ઘૂંટી} \gu{સૂજે} \gu{છે}?} — \emph{śũ tamārī ghū̃ṭī sūje che?} — \textquotedblleft{}Is your ankle swelling?\textquotedblright{}
\end{quote}

Not \textbf{\gu{ઘૂંટણ}}, the knee.

\begin{cousinweb}
\textbf{\gu{ઘૂંટી}} (\emph{ghū̃ṭī}) is inherited. CDIAL lists it, glossed \textquotedblleft{}ankle joint\textquotedblright{}, under the lexical Sanskrit \emph{ghuṭa}, \textquotedblleft{}ankle\textquotedblright{}, which it takes from Dravidian, after T. Burrow (CDIAL 4479).
\end{cousinweb}

\begin{grammarlens}[title={What you've built}]
One more word for everyday life at A2.
\end{grammarlens}

\subsection*{Your turn}
\begin{itemize}
\raggedright
  \item \emph{Say it:} \emph{ghū̃ṭī}
  \item \emph{Say it:} \emph{ghū̃ṭī}, once more
  \item \emph{Say it:} \emph{śũ tamārī ghū̃ṭī sūje che?} — Is your ankle swelling?
  \item \emph{Recall:} say the Gujarati for wrist, then the Gujarati for thumb, then say \emph{ghū̃ṭī} again
\end{itemize}

\begin{tcolorbox}[breakable,colback=teal!4,colframe=teal!35!black,title={Before you move on}]
What does \gu{ઘૂંટી} mean? (\textquotedblleft{}Ankle\textquotedblright{}.) Say it once more.
\end{tcolorbox}
\section[\texorpdfstring{\gu{એડી} (eḍī)}{એડી (eḍī)}]{\gu{એડી} (eḍī) — heel}

\label{lesson:GU-C229-edi}



\begin{quote}
Before the new one: say the Gujarati for thumb, then the Gujarati for ankle.
\end{quote}

\subsection*{You'll want to know: \gu{એડી}}
\textbf{\gu{એડી}} — \emph{eḍī} — \textquotedblleft{}heel\textquotedblright{}.

\begin{quote}
\textbf{\gu{મારી} \gu{એડી} \gu{નથી} \gu{દુખતી}.} — \emph{mārī eḍī nathī dukhtī} — \textquotedblleft{}My heel does not hurt.\textquotedblright{}
\end{quote}

\begin{cousinweb}
\textbf{\gu{એડી}} (\emph{eḍī}) comes from \emph{*aḍḍi}, \textquotedblleft{}heel\textquotedblright{}, a reconstructed form; CDIAL notes J. Bloch's view that it was borrowed from Dravidian, as in Tamil \emph{aḍi} (CDIAL 191).
\end{cousinweb}

\begin{grammarlens}[title={What you've built}]
One more word for everyday life at A2.
\end{grammarlens}

\subsection*{Your turn}
\begin{itemize}
\raggedright
  \item \emph{Say it:} \emph{eḍī}
  \item \emph{Say it:} \emph{eḍī}, once more
  \item \emph{Say it:} \emph{mārī eḍī nathī dukhtī} — My heel does not hurt.
  \item \emph{Recall:} say the Gujarati for thumb, then the Gujarati for ankle, then say \emph{eḍī} again
\end{itemize}

\begin{tcolorbox}[breakable,colback=teal!4,colframe=teal!35!black,title={Before you move on}]
What does \gu{એડી} mean? (\textquotedblleft{}Heel\textquotedblright{}.) Say it once more.
\end{tcolorbox}
\section[\texorpdfstring{\gu{પાંસળી} (pā̃sḷī)}{પાંસળી (pā̃sḷī)}]{\gu{પાંસળી} (pā̃sḷī) — rib}

\label{lesson:GU-C229-pasli}



\begin{quote}
Before the new one: say the Gujarati for ankle, then the Gujarati for heel.
\end{quote}

\subsection*{You'll want to know: \gu{પાંસળી}}
\textbf{\gu{પાંસળી}} — \emph{pā̃sḷī} — \textquotedblleft{}rib\textquotedblright{}.

\begin{quote}
\textbf{\gu{તેની} \gu{પાંસળીમાં} \gu{દુખાવો} \gu{નથી}.} — \emph{tenī pā̃sḷīmā̃ dukhāvo nathī} — \textquotedblleft{}There is no pain in his rib.\textquotedblright{}
\end{quote}

\begin{cousinweb}
\textbf{\gu{પાંસળી}} (\emph{pā̃sḷī}) is a tadbhava, a word inherited from Sanskrit and reshaped by sound change. It comes from \emph{parśu}, \textquotedblleft{}rib\textquotedblright{} (CDIAL 7948).
\end{cousinweb}

\begin{grammarlens}[title={What you've built}]
One more word for everyday life at A2.
\end{grammarlens}

\subsection*{Your turn}
\begin{itemize}
\raggedright
  \item \emph{Say it:} \emph{pā̃sḷī}
  \item \emph{Say it:} \emph{pā̃sḷī}, once more
  \item \emph{Say it:} \emph{tenī pā̃sḷīmā̃ dukhāvo nathī} — There is no pain in his rib.
  \item \emph{Recall:} say the Gujarati for ankle, then the Gujarati for heel, then say \emph{pā̃sḷī} again
\end{itemize}

\begin{tcolorbox}[breakable,colback=teal!4,colframe=teal!35!black,title={Before you move on}]
What does \gu{પાંસળી} mean? (\textquotedblleft{}Rib\textquotedblright{}.) Say it once more.
\end{tcolorbox}
